%!TeX root=../princesstop.tex
\chapter{The Indian Gentleman}

\begin{figure}[t!]
\centering
\includegraphics[width=\linewidth]{filigree}
\end{figure}

\lettrine[lines=5]{B}{ut} it was a perilous thing for Ermengarde and Lottie to make pilgrimages to the attic. They could never be quite sure when Sara would be there, and they could scarcely ever be certain that Miss Amelia would not make a tour of inspection through the bedrooms after the pupils were supposed to be asleep. So their visits were rare ones, and Sara lived a strange and lonely life. It was a lonelier life when she was downstairs than when she was in her attic. She had no one to talk to; and when she was sent out on errands and walked through the streets, a forlorn little figure carrying a basket or a parcel, trying to hold her hat on when the wind was blowing, and feeling the water soak through her shoes when it was raining, she felt as if the crowds hurrying past her made her loneliness greater. When she had been the Princess Sara, driving through the streets in her brougham, or walking, attended by Mariette, the sight of her bright, eager little face and picturesque coats and hats had often caused people to look after her. A happy, beautifully cared for little girl naturally attracts attention. Shabby, poorly dressed children are not rare enough and pretty enough to make people turn around to look at them and smile. No one looked at Sara in these days, and no one seemed to see her as she hurried along the crowded pavements. She had begun to grow very fast, and, as she was dressed only in such clothes as the plainer remnants of her wardrobe would supply, she knew she looked very queer, indeed. All her valuable garments had been disposed of, and such as had been left for her use she was expected to wear so long as she could put them on at all. Sometimes, when she passed a shop window with a mirror in it, she almost laughed outright on catching a glimpse of herself, and sometimes her face went red and she bit her lip and turned away.

In the evening, when she passed houses whose windows were lighted up, she used to look into the warm rooms and amuse herself by imagining things about the people she saw sitting before the fires or about the tables. It always interested her to catch glimpses of rooms before the shutters were closed. There were several families in the square in which Miss Minchin lived, with which she had become quite familiar in a way of her own. The one she liked best she called the Large Family. She called it the Large Family not because the members of it were big—for, indeed, most of them were little—but because there were so many of them. There were eight children in the Large Family, and a stout, rosy mother, and a stout, rosy father, and a stout, rosy grandmother, and any number of servants. The eight children were always either being taken out to walk or to ride in perambulators by comfortable nurses, or they were going to drive with their mamma, or they were flying to the door in the evening to meet their papa and kiss him and dance around him and drag off his overcoat and look in the pockets for packages, or they were crowding about the nursery windows and looking out and pushing each other and laughing—in fact, they were always doing something enjoyable and suited to the tastes of a large family. Sara was quite fond of them, and had given them names out of books—quite romantic names. She called them the Montmorencys when she did not call them the Large Family. The fat, fair baby with the lace cap was Ethelberta Beauchamp Montmorency; the next baby was Violet Cholmondeley Montmorency; the little boy who could just stagger and who had such round legs was Sydney Cecil Vivian Montmorency; and then came Lilian Evangeline Maud Marion, Rosalind Gladys, Guy Clarence, Veronica Eustacia, and Claude Harold Hector.

One evening a very funny thing happened—though, perhaps, in one sense it was not a funny thing at all.

Several of the Montmorencys were evidently going to a children's party, and just as Sara was about to pass the door they were crossing the pavement to get into the carriage which was waiting for them. Veronica Eustacia and Rosalind Gladys, in white-lace frocks and lovely sashes, had just got in, and Guy Clarence, aged five, was following them. He was such a pretty fellow and had such rosy cheeks and blue eyes, and such a darling little round head covered with curls, that Sara forgot her basket and shabby cloak altogether—in fact, forgot everything but that she wanted to look at him for a moment. So she paused and looked.

It was Christmas time, and the Large Family had been hearing many stories about children who were poor and had no mammas and papas to fill their stockings and take them to the pantomime—children who were, in fact, cold and thinly clad and hungry. In the stories, kind people—sometimes little boys and girls with tender hearts—invariably saw the poor children and gave them money or rich gifts, or took them home to beautiful dinners. Guy Clarence had been affected to tears that very afternoon by the reading of such a story, and he had burned with a desire to find such a poor child and give her a certain sixpence he possessed, and thus provide for her for life. An entire sixpence, he was sure, would mean affluence for evermore. As he crossed the strip of red carpet laid across the pavement from the door to the carriage, he had this very sixpence in the pocket of his very short man-o-war trousers; And just as Rosalind Gladys got into the vehicle and jumped on the seat in order to feel the cushions spring under her, he saw Sara standing on the wet pavement in her shabby frock and hat, with her old basket on her arm, looking at him hungrily.

He thought that her eyes looked hungry because she had perhaps had nothing to eat for a long time. He did not know that they looked so because she was hungry for the warm, merry life his home held and his rosy face spoke of, and that she had a hungry wish to snatch him in her arms and kiss him. He only knew that she had big eyes and a thin face and thin legs and a common basket and poor clothes. So he put his hand in his pocket and found his sixpence and walked up to her benignly.

<Here, poor little girl,> he said. <Here is a sixpence. I will give it to you.>

Sara started, and all at once realized that she looked exactly like poor children she had seen, in her better days, waiting on the pavement to watch her as she got out of her brougham. And she had given them pennies many a time. Her face went red and then it went pale, and for a second she felt as if she could not take the dear little sixpence.

<Oh, no!> she said. <Oh, no, thank you; I mustn't take it, indeed!>

Her voice was so unlike an ordinary street child's voice and her manner was so like the manner of a well-bred little person that Veronica Eustacia (whose real name was Janet) and Rosalind Gladys (who was really called Nora) leaned forward to listen.

But Guy Clarence was not to be thwarted in his benevolence. He thrust the sixpence into her hand.

<Yes, you must take it, poor little girl!> he insisted stoutly. <You can buy things to eat with it. It is a whole sixpence!>

There was something so honest and kind in his face, and he looked so likely to be heartbrokenly disappointed if she did not take it, that Sara knew she must not refuse him. To be as proud as that would be a cruel thing. So she actually put her pride in her pocket, though it must be admitted her cheeks burned.

<Thank you,> she said. <You are a kind, kind little darling thing.> And as he scrambled joyfully into the carriage she went away, trying to smile, though she caught her breath quickly and her eyes were shining through a mist. She had known that she looked odd and shabby, but until now she had not known that she might be taken for a beggar.

As the Large Family's carriage drove away, the children inside it were talking with interested excitement.

<Oh, Donald,> (this was Guy Clarence's name), Janet exclaimed alarmedly, <why did you offer that little girl your sixpence? I'm sure she is not a beggar!>

<She didn't speak like a beggar!> cried Nora. <And her face didn't really look like a beggar's face!>

<Besides, she didn't beg,> said Janet. <I was so afraid she might be angry with you. You know, it makes people angry to be taken for beggars when they are not beggars.>

<She wasn't angry,> said Donald, a trifle dismayed, but still firm. <She laughed a little, and she said I was a kind, kind little darling thing. And I was!>—stoutly. <It was my whole sixpence.>

Janet and Nora exchanged glances.

<A beggar girl would never have said that,> decided Janet. <She would have said, <Thank yer kindly, little gentleman—thank yer, sir;> and perhaps she would have bobbed a curtsy.>

Sara knew nothing about the fact, but from that time the Large Family was as profoundly interested in her as she was in it. Faces used to appear at the nursery windows when she passed, and many discussions concerning her were held round the fire.

<She is a kind of servant at the seminary,> Janet said. <I don't believe she belongs to anybody. I believe she is an orphan. But she is not a beggar, however shabby she looks.>

And afterward she was called by all of them, <The-little-girl-who-is-not-a-beggar,> which was, of course, rather a long name, and sounded very funny sometimes when the youngest ones said it in a hurry.

Sara managed to bore a hole in the sixpence and hung it on an old bit of narrow ribbon round her neck. Her affection for the Large Family increased—as, indeed, her affection for everything she could love increased. She grew fonder and fonder of Becky, and she used to look forward to the two mornings a week when she went into the schoolroom to give the little ones their French lesson. Her small pupils loved her, and strove with each other for the privilege of standing close to her and insinuating their small hands into hers. It fed her hungry heart to feel them nestling up to her. She made such friends with the sparrows that when she stood upon the table, put her head and shoulders out of the attic window, and chirped, she heard almost immediately a flutter of wings and answering twitters, and a little flock of dingy town birds appeared and alighted on the slates to talk to her and make much of the crumbs she scattered. With Melchisedec she had become so intimate that he actually brought Mrs Melchisedec with him sometimes, and now and then one or two of his children. She used to talk to him, and, somehow, he looked quite as if he understood.

There had grown in her mind rather a strange feeling about Emily, who always sat and looked on at everything. It arose in one of her moments of great desolateness. She would have liked to believe or pretend to believe that Emily understood and sympathized with her. She did not like to own to herself that her only companion could feel and hear nothing. She used to put her in a chair sometimes and sit opposite to her on the old red footstool, and stare and pretend about her until her own eyes would grow large with something which was almost like fear—particularly at night when everything was so still, when the only sound in the attic was the occasional sudden scurry and squeak of Melchisedec's family in the wall. One of her <pretends> was that Emily was a kind of good witch who could protect her. Sometimes, after she had stared at her until she was wrought up to the highest pitch of fancifulness, she would ask her questions and find herself \textsc{almost} feeling as if she would presently answer. But she never did.

<As to answering, though,> said Sara, trying to console herself, <I don't answer very often. I never answer when I can help it. When people are insulting you, there is nothing so good for them as not to say a word—just to look at them and \textsc{think}. Miss Minchin turns pale with rage when I do it, Miss Amelia looks frightened, and so do the girls. When you will not fly into a passion people know you are stronger than they are, because you are strong enough to hold in your rage, and they are not, and they say stupid things they wish they hadn't said afterward. There's nothing so strong as rage, except what makes you hold it in—that's stronger. It's a good thing not to answer your enemies. I scarcely ever do. Perhaps Emily is more like me than I am like myself. Perhaps she would rather not answer her friends, even. She keeps it all in her heart.>

But though she tried to satisfy herself with these arguments, she did not find it easy. When, after a long, hard day, in which she had been sent here and there, sometimes on long errands through wind and cold and rain, she came in wet and hungry, and was sent out again because nobody chose to remember that she was only a child, and that her slim legs might be tired and her small body might be chilled; when she had been given only harsh words and cold, slighting looks for thanks; when the cook had been vulgar and insolent; when Miss Minchin had been in her worst mood, and when she had seen the girls sneering among themselves at her shabbiness—then she was not always able to comfort her sore, proud, desolate heart with fancies when Emily merely sat upright in her old chair and stared.

One of these nights, when she came up to the attic cold and hungry, with a tempest raging in her young breast, Emily's stare seemed so vacant, her sawdust legs and arms so inexpressive, that Sara lost all control over herself. There was nobody but Emily—no one in the world. And there she sat.

<I shall die presently,> she said at first.

Emily simply stared.

<I can't bear this,> said the poor child, trembling. <I know I shall die. I'm cold; I'm wet; I'm starving to death. I've walked a thousand miles today, and they have done nothing but scold me from morning until night. And because I could not find that last thing the cook sent me for, they would not give me any supper. Some men laughed at me because my old shoes made me slip down in the mud. I'm covered with mud now. And they laughed. Do you hear?>

She looked at the staring glass eyes and complacent face, and suddenly a sort of heartbroken rage seized her. She lifted her little savage hand and knocked Emily off the chair, bursting into a passion of sobbing—Sara who never cried.

<You are nothing but a \textsc{doll}!> she cried. <Nothing but a doll—doll—doll! You care for nothing. You are stuffed with sawdust. You never had a heart. Nothing could ever make you feel. You are a \textsc{doll}!> Emily lay on the floor, with her legs ignominiously doubled up over her head, and a new flat place on the end of her nose; but she was calm, even dignified. Sara hid her face in her arms. The rats in the wall began to fight and bite each other and squeak and scramble. Melchisedec was chastising some of his family.

Sara's sobs gradually quieted themselves. It was so unlike her to break down that she was surprised at herself. After a while she raised her face and looked at Emily, who seemed to be gazing at her round the side of one angle, and, somehow, by this time actually with a kind of glassy-eyed sympathy. Sara bent and picked her up. Remorse overtook her. She even smiled at herself a very little smile.

<You can't help being a doll,> she said with a resigned sigh, <any more than Lavinia and Jessie can help not having any sense. We are not all made alike. Perhaps you do your sawdust best.> And she kissed her and shook her clothes straight, and put her back upon her chair.

She had wished very much that some one would take the empty house next door. She wished it because of the attic window which was so near hers. It seemed as if it would be so nice to see it propped open someday and a head and shoulders rising out of the square aperture.

<If it looked a nice head,> she thought, <I might begin by saying, <Good morning,> and all sorts of things might happen. But, of course, it's not really likely that anyone but under servants would sleep there.>

One morning, on turning the corner of the square after a visit to the grocer's, the butcher's, and the baker's, she saw, to her great delight, that during her rather prolonged absence, a van full of furniture had stopped before the next house, the front doors were thrown open, and men in shirt sleeves were going in and out carrying heavy packages and pieces of furniture.

<It's taken!> she said. <It really \textsc{is} taken! Oh, I do hope a nice head will look out of the attic window!>

She would almost have liked to join the group of loiterers who had stopped on the pavement to watch the things carried in. She had an idea that if she could see some of the furniture she could guess something about the people it belonged to.

<Miss Minchin's tables and chairs are just like her,> she thought; <I remember thinking that the first minute I saw her, even though I was so little. I told papa afterward, and he laughed and said it was true. I am sure the Large Family have fat, comfortable armchairs and sofas, and I can see that their red-flowery wallpaper is exactly like them. It's warm and cheerful and kind-looking and happy.>

She was sent out for parsley to the greengrocer's later in the day, and when she came up the area steps her heart gave quite a quick beat of recognition. Several pieces of furniture had been set out of the van upon the pavement. There was a beautiful table of elaborately wrought teakwood, and some chairs, and a screen covered with rich Oriental embroidery. The sight of them gave her a weird, homesick feeling. She had seen things so like them in India. One of the things Miss Minchin had taken from her was a carved teakwood desk her father had sent her.

<They are beautiful things,> she said; <they look as if they ought to belong to a nice person. All the things look rather grand. I suppose it is a rich family.>

The vans of furniture came and were unloaded and gave place to others all the day. Several times it so happened that Sara had an opportunity of seeing things carried in. It became plain that she had been right in guessing that the newcomers were people of large means. All the furniture was rich and beautiful, and a great deal of it was Oriental. Wonderful rugs and draperies and ornaments were taken from the vans, many pictures, and books enough for a library. Among other things there was a superb god Buddha in a splendid shrine.

<Someone in the family \textsc{must} have been in India,> Sara thought. <They have got used to Indian things and like them. I \textsc{am} glad. I shall feel as if they were friends, even if a head never looks out of the attic window.>

When she was taking in the evening's milk for the cook (there was really no odd job she was not called upon to do), she saw something occur which made the situation more interesting than ever. The handsome, rosy man who was the father of the Large Family walked across the square in the most matter-of-fact manner, and ran up the steps of the next-door house. He ran up them as if he felt quite at home and expected to run up and down them many a time in the future. He stayed inside quite a long time, and several times came out and gave directions to the workmen, as if he had a right to do so. It was quite certain that he was in some intimate way connected with the newcomers and was acting for them.

<If the new people have children,> Sara speculated, <the Large Family children will be sure to come and play with them, and they \textsc{might} come up into the attic just for fun.>

At night, after her work was done, Becky came in to see her fellow prisoner and bring her news.

<It's a' Nindian gentleman that's comin' to live next door, miss,> she said. <I don't know whether he's a black gentleman or not, but he's a Nindian one. He's very rich, an' he's ill, an' the gentleman of the Large Family is his lawyer. He's had a lot of trouble, an' it's made him ill an' low in his mind. He worships idols, miss. He's an 'eathen an' bows down to wood an' stone. I seen a' idol bein' carried in for him to worship. Somebody had oughter send him a trac'. You can get a trac' for a penny.>

Sara laughed a little.

<I don't believe he worships that idol,> she said; <some people like to keep them to look at because they are interesting. My papa had a beautiful one, and he did not worship it.>

But Becky was rather inclined to prefer to believe that the new neighbour was <an 'eathen.> It sounded so much more romantic than that he should merely be the ordinary kind of gentleman who went to church with a prayer book. She sat and talked long that night of what he would be like, of what his wife would be like if he had one, and of what his children would be like if they had children. Sara saw that privately she could not help hoping very much that they would all be black, and would wear turbans, and, above all, that—like their parent—they would all be <'eathens.>

<I never lived next door to no 'eathens, miss,> she said; <I should like to see what sort o' ways they'd have.>

It was several weeks before her curiosity was satisfied, and then it was revealed that the new occupant had neither wife nor children. He was a solitary man with no family at all, and it was evident that he was shattered in health and unhappy in mind.

A carriage drove up one day and stopped before the house. When the footman dismounted from the box and opened the door the gentleman who was the father of the Large Family got out first. After him there descended a nurse in uniform, then came down the steps two men-servants. They came to assist their master, who, when he was helped out of the carriage, proved to be a man with a haggard, distressed face, and a skeleton body wrapped in furs. He was carried up the steps, and the head of the Large Family went with him, looking very anxious. Shortly afterward a doctor's carriage arrived, and the doctor went in—plainly to take care of him.

<There is such a yellow gentleman next door, Sara,> Lottie whispered at the French class afterward. <Do you think he is a Chinee? The geography says the Chinee men are yellow.>

<No, he is not Chinese,> Sara whispered back; <he is very ill. Go on with your exercise, Lottie. <\textfrench{Non, monsieur. Je n'ai pas le canif de mon oncle.}>>

That was the beginning of the story of the Indian gentleman.